\section{Comandi per l\textquotesingle Architettura di
Rete}\label{comandi-per-larchitettura-di-rete}

Per andare a fondo nei dettagli implementativi e spiegare meglio alcuni
concetti, sono riportati alcuni passaggi fondamentali nella
configurazione dell\textquotesingle host che esegue BlueAgent.

\subsection{configurare OVN}\label{configurare-ovn}

Qui riporto la configurazione per implementare la SDN. Il \textbf{piano
dei dati} e\textquotesingle{} implementato da Open vSwitch
(\textbf{OVS}) e deve esser configurato in modo che possa comunicare con
Incus.

\begin{quote}
\passthrough{\lstinline!SDN (Software Defined Networking)!}: Paradigma
architetturale che separa il \textbf{Control Plane} (logica decisionale
centralizzata) dal \textbf{Data Plane} (inoltro fisico dei pacchetti),
consentendo di gestire e programmare l\textquotesingle infrastruttura di
rete interamente via software.

\passthrough{\lstinline!Open vSwitch (OVS)!}: Switch virtuale operante
nel kernel Linux come Data Plane. Sostituisce il classico bridge Linux
eseguendo l\textquotesingle inoltro dei frame ad alte prestazioni
basandosi su tabelle di flussi (\emph{flow tables}).

\passthrough{\lstinline!OpenFlow!}: Protocollo standard che permette al
controller SDN di comunicare con lo switch e programmare le regole di
matching e le azioni direttamente all\textquotesingle interno delle
pipeline di Open vSwitch.

\passthrough{\lstinline!OVN (Open Virtual Network)!}: Sistema SDN che
opera come Control Plane per Open vSwitch. Traduce topologie logiche
(Logical Switch, Logical Router, NAT e ACL) in flussi OpenFlow
distribuiti, istruendo Open vSwitch su come instradare il traffico senza
creare molteplici interfacce o bridge dedicati nel kernel host.
\end{quote}

\textbf{Configurazione di open vSwitch}:

\begin{lstlisting}[language=bash]
sudo ovs-vsctl set open_vswitch . \ 
external_ids:ovn-remote=unix:/run/ovn/ovnsb_db.sock \ external_ids:ovn-encap-type=geneve \ 
external_ids:ovn-encap-ip=127.0.0.1
\end{lstlisting}

\begin{itemize}
\tightlist
\item
  \textbf{\passthrough{\lstinline!ovn-encap-type=geneve!}}:
  specifichiamo geneve come protocollo da utilizzare per il tunneling.
\item
  \textbf{\passthrough{\lstinline!ovn-encap-ip=127.0.0.1!}}: imposta la
  terminazione del tunnel nel localhost.
\end{itemize}

\begin{quote}
\textbf{Open vSwitch}: e\textquotesingle{} uno switch virtuale,
software-definied. Agisce per la maggior parte in user space e comunica
con il kernel per riceve i pacchetti ed inoltrarli.

\begin{itemize}
\tightlist
\item
  \passthrough{\lstinline!ovs-vswitchd!}: e\textquotesingle{} il demone
  in user space. Ispeziona i pacchetti catturati dal kernel.
\item
  \passthrough{\lstinline!ovsdb-server!}: e\textquotesingle{} un
  database che memorizza la configurazione degli switch
\item
  \passthrough{\lstinline!ovs-vsctl!}: utility da riga di comando per la
  gestione dell\textquotesingle infrastruttura.
\end{itemize}
\end{quote}

\textbf{Collegare l\textquotesingle interfaccia Northbound di OV ad
Incus}:

\begin{lstlisting}[language=bash]
ovn-nbctl set-connection ptcp:6641:127.0.0.1
incus config set network.ovn.northbound_connection tcp:127.0.0.1:6641
\end{lstlisting}

\begin{quote}
\begin{itemize}
\tightlist
\item
  \textbf{OVN Northbound}: e\textquotesingle{}
  l\textquotesingle interfaccia dove dichiaro le reti che voglio creare
  e le loro proprietà che devono rispettare.
\item
  \textbf{\passthrough{\lstinline!ovn-northd!}}: il demone che traduce
  la topologia dichiarata tramite l\textquotesingle interfaccia north
  bound in stati logici.
\item
  \textbf{OVN Southbound (SB) DB:}~Memorizza le tabelle di flussi logici
  compilate e le associazioni di chassis (host fisici)
\end{itemize}

\begin{itemize}
\tightlist
\item
  \textbf{\passthrough{\lstinline!ovn-controller!}:}~Un agente in
  esecuzione sull\textquotesingle host che legge il Southbound DB e
  programma direttamente Open vSwitch locale.
\end{itemize}
\end{quote}

\subsection{\texorpdfstring{\(\text{production-}k\) Routing e
Connettività dei Peer (WireGuard e
nftables)}{\textbackslash text\{production-\}k Routing e Connettività dei Peer (WireGuard e nftables)}}\label{textproduction-k-routing-e-connettivituxe0-dei-peer-wireguard-e-nftables}

L\textquotesingle accesso dei partecipanti ai container rete
\(\text{production-}k\) richiede l\textquotesingle applicazione di
regole di inoltro esplicite tra \passthrough{\lstinline!wg-k!} ed
\passthrough{\lstinline!br-team-k!}, da applicare nel kernel Linux
mediante \passthrough{\lstinline!nftables!}.

\subsubsection{\texorpdfstring{Inoltro Kernel e Controllo di Flusso
(\texttt{nftables})}{Inoltro Kernel e Controllo di Flusso (nftables)}}\label{inoltro-kernel-e-controllo-di-flusso-nftables}

Le regole per il traffico vengono definite dentro la tabella
\passthrough{\lstinline!inet blueagent!}:

\begin{lstlisting}[language=bash]
sudo nft add table inet blueagent
sudo nft add set inet blueagent team_links '{ type ifname . ifname; flags interval; }'

sudo nft add element inet blueagent team_links { "wg-$k" . "br-team-$k", "br-team-$k" . "wg-$k" }
\end{lstlisting}

All\textquotesingle interno della catena
\passthrough{\lstinline!forward!}, il traffico viene validato
verificando l\textquotesingle appartenenza al set
\passthrough{\lstinline!@team\_links!} e scartando esplicitamente le
comunicazioni incrociate tra team diversi. Vengono inserite
esplicitamente per ogni \passthrough{\lstinline!k!} le coppie di
interfacce per cui autorizzare il traffico.

A priori, si vuole impedire l\textquotesingle inoltro dei pacchetti
applicando l\textquotesingle azione \passthrough{\lstinline!drop!} tra
le interfacce \passthrough{\lstinline!wg!} e
\passthrough{\lstinline!br-team!}, in modo da impedire rotte che non
sono state esplicitamente accettate nella tabella:

\begin{lstlisting}[language=bash]
sudo nft add rule inet blueagent forward iifname "wg-*" oifname "br-team-*" drop
sudo nft add rule inet blueagent forward iifname "br-team-*" oifname "wg-*" drop
sudo nft add rule inet blueagent forward iifname "br-team-*" oifname "br-team-*" drop
sudo nft add rule inet blueagent forward iifname "wg-*" oifname "wg-*" drop
\end{lstlisting}

\subsubsection{Connettività Egress
(Masquerading)}\label{connettivituxe0-egress-masquerading}

Per consentire l\textquotesingle accesso a Internet ai container e ai
peer, viene applicato il mascheramento dell\textquotesingle indirizzo
sorgente verso l\textquotesingle esterno dell\textquotesingle envelope.
Si abilita il \passthrough{\lstinline!nat!} (dunque, il meccanismo di
traduzione degli indirizzi) poco prima che il frame esca e venga immesso
fisicamente sulla porta. Poi specifichiamo che si applica il nat
dinamico (\passthrough{\lstinline!masquerade!}) per gli indirizzi che
provengono dalla rete dei container verso l\textquotesingle esterno.

\begin{lstlisting}[language=bash]
sudo nft add chain inet blueagent postrouting '{ type nat hook postrouting priority srcnat; policy accept; }'
sudo nft add rule inet blueagent postrouting ip saddr 10.0.0.0/8 ip daddr != 10.0.0.0/8 masquerade
\end{lstlisting}
